% ===================================================================
%  CADRO N.º 15 — O mercado. — Os comestibles.
%  Traduction 2026 d'apres texto/io, controlee sur texto/fr, texto/es
%  et texto/pt. Consigne : voir texto/gl/10-cadro-01.tex.
%
%  LE TABLEAU DU MARCHE EST CELUI OU LA REGLE DU LIEU D'APPRENTISSAGE
%  SE VERIFIE LE MIEUX DE TOUTE LA COLONNE, comme elle s'etait
%  verifiee au meme tableau en irlandais et pour la meme raison : la
%  nourriture s'apprend a table, c'est-a-dire chez soi. Cent objets,
%  et pas un emprunt parmi eux : « ovos », « leituga », « repolo »,
%  « cogomelos », « ameixas », « abelás », « touciño », « chourizo »,
%  « fabas », « améndoas », « noces », « pexegos ».
%
%  ET C'EST L'ETAL DU POISSONNIER QUI LE MONTRE LE PLUS NETTEMENT.
%  Onze mots d'affilee, sans un seul emprunt : « anguías »,
%  « linguados », « troitas », « carpas », « bacallau »,
%  « mexillóns », « atún », « ostras », « camaróns », « cangrexos »,
%  « lumbrigantes ». C'est la regle du tableau 8 verifiee une seconde
%  fois : LA OU L'ON DISTINGUE, ON A DES MOTS — et l'on distingue la
%  ou l'on vit de la chose. La vendange en avait donne quatre pour un
%  tonneau ; la criee en donne onze pour un etal.
% ===================================================================

%%K t15-apar-1 apar
\VUsaut{4.0mm}
\VUtitre{15.0pt}{60}{CADRO N.\textsuperscript{o} 15}
\VUsaut{3.0mm}

%%K t15-tit-1 sub
\VUpk{11.4pt}{40}{O mercado. --- Os comestibles.}
\VUsaut{3.0mm}

%%K t15-01-1 p
1. --- A maior parte dos comerciantes que nos fornecen as mercadorías precisas para o
noso sustento xúntanse baixo a \VUgras{praza cuberta} \textsuperscript{(1)} do
\VUgras{mercado} ou nas rúas veciñas.

%%K t15-02-1 p
2. --- O \VUgras{chacineiro} \textsuperscript{(2)}, que vende \VUgras{xamón}
\textsuperscript{(3)}, \VUgras{morcilla} \textsuperscript{(4)}, \VUgras{chourizo}
\textsuperscript{(5)}, \VUgras{androlla} \textsuperscript{(6)} e \VUgras{touciño}
\textsuperscript{(7)}, está preto da \VUgras{vendedora de manteiga e queixo}
\textsuperscript{(8)}, cuxo posto está provisto de \VUgras{queixos holandeses}
\textsuperscript{(9)} de codia vermella, de \VUgras{queixos camembert}
\textsuperscript{(10)}, de \VUgras{rodas de gruyère} \textsuperscript{(11)} e de
\VUgras{pelotas de manteiga} fresca \textsuperscript{(12)}.

%%K t15-03-1 p
3. --- Diante dela, unha \VUgras{vendedora de legumes} \textsuperscript{(13)}
ofrécelles ás súas freguesas belos \VUgras{tomates} \textsuperscript{(14)},
\VUgras{coliflores} \textsuperscript{(15)}, \VUgras{leituga}
\textsuperscript{(16)}, \VUgras{repolos} \textsuperscript{(17)},
\VUgras{cabaciñas} \textsuperscript{(19)}, \VUgras{melóns} \textsuperscript{(18)} e
mesmo \VUgras{cogomelos} \textsuperscript{(20)}. Pero un \VUgras{mozo da
cervexaría}, que parou o seu \VUgras{carro de reparto} \textsuperscript{(21)}
cargado de \VUgras{barrís de cervexa} \textsuperscript{(22)} xusto diante do posto
dela, impide que as súas freguesas se acheguen. Unha horticultora tráelle un cesto
grande de \VUgras{alcachofas} \textsuperscript{(23)}.

%%K t15-tit-2 sub
\VUsaut{4.0mm}
\VUpk{10.2pt}{40}{Venda e compra.}
\VUsaut{3.0mm}

%%K t15-04-1 p
4. --- No mercado hai moita animación, pois chegou a hora da \VUgras{poxa}
\textsuperscript{(24)}. As \VUgras{donas da casa} \textsuperscript{(26)}, que veñen
alí facer as súas compras, paran ao pasar diante da tenda da \VUgras{tripeira}
\textsuperscript{(25)} para mercaren \VUgras{callos} ou unha \VUgras{cabeza de
xato.}

%%K t15-05-1 p
5. --- Un \VUgras{cociñeiro} \textsuperscript{(27)} graxento regatea un \VUgras{atún}
moi belo á \VUgras{peixeira} \textsuperscript{(28)}, que sobe os seus prezos ao seu
antollo. Recibiu unha gran cantidade de \VUgras{anguías, linguados, troitas} e
\VUgras{carpas.} O seu \VUgras{bacallau} e os seus \VUgras{mexillóns} están moi
frescos.

%%K t15-tit-3 sub
\VUsaut{4.0mm}
\VUpk{10.2pt}{40}{Mercadorías baratas e caras.}
\VUsaut{3.0mm}

%%K t15-06-1 p
6. --- A \VUgras{vendedora de aves} \textsuperscript{(29)}, instalada preto dela, é
moi honrada. Cando vende un \VUgras{pavo,} unha \VUgras{oca,} un \VUgras{pato} ou un
simple \VUgras{polo,} só pide o prezo xusto.

%%K t15-07-1 p
7. --- Na esquina da praza, no primeiro andar, hai instalado desde hai pouco tempo un
\VUgras{vendedor de tea} \textsuperscript{(30)} na casa do cal as compradoras están
sempre seguras de atopar unha moi bela variedade de \VUgras{manteis,}
\VUgras{gardanapos, mandís} e \VUgras{panos de cociña.} No mesmo andar, un
\VUgras{coiteleiro} \textsuperscript{(31)} instalou o seu taller. Afía os
\VUgras{coitelos} das vendedoras da praza e fabrica para os horteláns \VUgras{fouciños}
do \VUgras{aceiro} máis duro.

%%K t15-tit-4 sub
\VUsaut{4.0mm}
\VUpk{10.2pt}{40}{A tenda de ultramarinos.}
\VUsaut{3.0mm}

%%K t15-08-1 p
8. --- Baixo o taller del abriuse hai pouco tempo un gran almacén de alimentos. Xa non
é a simple e pouco importante \VUgras{tenda de especias} \textsuperscript{(32)} de
antano, que só vendía as mercadorías correntes, \VUgras{té} \textsuperscript{(33)},
\VUgras{chocolate} \textsuperscript{(34)}, \VUgras{azucre en pan}
\textsuperscript{(37)} e \VUgras{xabón} \textsuperscript{(38)}. Alí véndese calquera
mercadoría, \VUgras{galletas crocantes}, \VUgras{conservas}
\textsuperscript{(35)}, \VUgras{licores} \textsuperscript{(36)}, \VUgras{ron,}
\VUgras{coñac, kirsch, anís} e \VUgras{curaçao.} Alí atópase caza: \VUgras{lebre}
\textsuperscript{(39)}, \VUgras{tordos} \textsuperscript{(40)}, \VUgras{perdices}
\textsuperscript{(41)}, \VUgras{codornices} \textsuperscript{(42)}, \VUgras{pato
bravo} \textsuperscript{(43)} ou \VUgras{faisán} \textsuperscript{(44)}, tan
doadamente coma froitas exóticas tales coma \VUgras{plátanos}
\textsuperscript{(46)} e \VUgras{ananases.}

%%K t15-09-1 p
9. --- Un \VUgras{mozo} de ultramarinos \textsuperscript{(45)} moi amable está
disposto a dar o que un queira: \VUgras{marmelada} \textsuperscript{(47)} de grosellas
ou de marmelos, \VUgras{graxa} \textsuperscript{(48)}, \VUgras{pepiniños}
\textsuperscript{(49)} ou \VUgras{sardiñas} \textsuperscript{(50)} salgadas.

%%K t15-09-2 p
Iso é moi cómodo para as \VUgras{freguesas} \textsuperscript{(51)}, que atopan, ao
lado das mercadorías máis correntes, as primicias máis raras e as froitas máis
saborosas: \VUgras{peras} \textsuperscript{(52)} zumarentas, \VUgras{ameixas}
\textsuperscript{(53)}, \VUgras{uvas} \textsuperscript{(54)}, \VUgras{pexegos}
\textsuperscript{(55)}, \VUgras{améndoas} \textsuperscript{(56)} e \VUgras{noces}
\textsuperscript{(57)}.

%%K t15-tit-5 sub
\VUsaut{4.0mm}
\VUpk{10.2pt}{40}{A clientela.}
\VUsaut{3.0mm}

%%K t15-10-1 p
10. --- O \VUgras{arroz} \textsuperscript{(58)} e as \VUgras{lentellas}
\textsuperscript{(59)} están ao lado da caixa dos \VUgras{arenques afumados}
\textsuperscript{(60)}; a \VUgras{pataca} \textsuperscript{(61)} e as \VUgras{fabas}
\textsuperscript{(63)} están xunto ás \VUgras{mazás} \textsuperscript{(62)}, aos
\VUgras{figos} \textsuperscript{(64)}, ás \VUgras{ameixas pasas}
\textsuperscript{(65)} e ás \VUgras{abelás} \textsuperscript{(66)}. Nese almacén
atópanse \VUgras{ovos} \textsuperscript{(67)} frescos e \VUgras{olivas}
\textsuperscript{(68)}; por iso os \VUgras{cestos} \textsuperscript{(70)} e as
\VUgras{redes da compra} \textsuperscript{(73)} se enchen axiña.

%%K t15-11-1 p
11. --- O \VUgras{café} \textsuperscript{(72)} desa excelente casa acadou, entre as
\VUgras{criadas} \textsuperscript{(69)} e as \VUgras{cociñeiras,} unha xusta sona,
pois o vello \VUgras{tendeiro} \textsuperscript{(71)} en persoa tórrao co maior
coidado.

%%K t15-tit-6 sub
\VUsaut{4.0mm}
\VUpk{10.2pt}{40}{Os espectáculos.}
\VUsaut{3.0mm}

%%K t15-12-1 p
12. --- Non todos van ao mercado para se abasteceren. No movemento pintoresco da rúa
pequena que leva á praza cuberta vense as persoas máis diversas. Ao lado dun amable
\VUgras{mozo de carnicería} \textsuperscript{(75)}, que empurra diante de si unha
\VUgras{carretilla} \textsuperscript{(74)}, velaquí dous soldados de permiso, moi
orgullosos de amosaren o seu uniforme brillante: o \VUgras{húsar}
\textsuperscript{(76)} co seu \VUgras{dolmán} azul e o seu \VUgras{chacó de pluma,} e
o \VUgras{cabo de zuavos,} coa manga inchada e a cabeza cuberta cun \VUgras{fez.}

%%K t15-13-1 p
13. --- O \VUgras{panadeiro} \textsuperscript{(80)}, que está de pé na soleira da
\VUgras{panadaría} \textsuperscript{(78)}, acaba de amasar a masa do seu \VUgras{pan}
\textsuperscript{(79)} e de sacar do forno os \VUgras{croissants}
\textsuperscript{(81)}, os \VUgras{bolos} \textsuperscript{(82)} e as \VUgras{fogazas}
\textsuperscript{(83)} que están no escaparate.

%%K t15-14-1 p
14. --- Por riba da panadaría vive unha hábil \VUgras{costureira de roupa branca}
\textsuperscript{(84)} que emprega varias \VUgras{bordadoras} \textsuperscript{(85)}.
Ao lado dela vive unha \VUgras{lavandeira e pasadora} \textsuperscript{(86)}, que
almidoa a \VUgras{roupa branca} \textsuperscript{(88)} despois de a lavar e secar, e
que a pasa co \VUgras{ferro} \textsuperscript{(87)}.

%%K t15-tit-7 sub
\VUsaut{4.0mm}
\VUpk{10.2pt}{40}{Carne, peixe e legumes.}
\VUsaut{3.0mm}

%%K t15-15-1 p
15. --- Na \VUgras{carnicería} \textsuperscript{(89)}, situada no andar baixo,
véndese toda clase de \VUgras{carnes, carne de carneiro} \textsuperscript{(90)},
\VUgras{carne de vaca} \textsuperscript{(94)} ou \VUgras{carne de xato}
\textsuperscript{(92)} en forma de \VUgras{pernas de carneiro, bistés, asados} e
\VUgras{costeletas.} O \VUgras{mozo do carniceiro} \textsuperscript{(93)} corta todas
esas carnes e pon aparte os \VUgras{ósos con tutano} \textsuperscript{(96)}. Á porta da
carnicería está unha \VUgras{ostreira} \textsuperscript{(97)} que vende
\VUgras{ostras} \textsuperscript{(98)}. Ábreas se un o desexa.

%%K t15-16-1 p
16. --- Todos eses comerciantes pagan unha patente ou unha taxa de posto; por iso o
\VUgras{policía} \textsuperscript{(100)} persegue sen piedade as pobres
\VUgras{vendedoras ambulantes de legumes} \textsuperscript{(99)}, que lles fan a
competencia levando nos seus carriños \VUgras{cenorias, rabanetes, nabos, porros} ou
\VUgras{mollos de espárragos, chícharos frescos, espinacas, acedeiras, agrón} e
\VUgras{perexil.}

%%K t15-tit-8 sub
\VUsaut{4.0mm}
\VUpk{10.2pt}{40}{Coidado co policía!}
\VUsaut{3.0mm}

%%K t15-17-1 p
17. --- O rapaz que vende \VUgras{allos} \textsuperscript{(101)} e \VUgras{cebolas}
sabe moi ben que tampouco el pode vender a súa mercadoría preto do mercado. Foxe,
correndo co risco de derrubar o cesto de \VUgras{cangrexos,} \VUgras{cangrexos de río} e
\VUgras{camaróns} do \VUgras{vendedor} \textsuperscript{(102)} de \VUgras{lagostas} e
\VUgras{lumbrigantes.}

%%K t15-18-1 p
18. --- Todos se axitan e fan moito ruído; pero iso non impide oír claramente a voz do
\VUgras{vidreiro} \textsuperscript{(103)} ambulante, nin o berro do \VUgras{carboeiro}
\textsuperscript{(104)} que acaba de deixar uns intres o seu carriño para entregar un
\VUgras{saco de carbón} \textsuperscript{(105)}.
